\documentclass[../../../../main.tex]{subfiles}
\begin{document}

{\bf [解]}
題意より$g(x) = \alpha x + \beta \quad (\alpha \ne 0, \, \alpha, \beta \in \mathbb{R})$ とおける．条件より
\begin{align}
\alpha \int_0^a f(x) x \, dx + \beta \int_0^a f(x) \, dx = 0 \label{eq:1}
\end{align}
が $\alpha \ne 0$ をみたす任意の $(\alpha, \beta)$ で成立するので\cref{eq:1}が $\alpha, \beta$ についての恒等式となる．従って
\begin{align}
\int_0^a f(x) \, dx &= 0 \label{eq:2} \\
\int_0^a f(x) x \, dx &= 0 \label{eq:3}
\end{align}
である．以下この条件を満たす$f(x)$について考える．

\paragraph{$f(x)=0$が実数解を持たない時}

$f(x) = 0$ が実解を持たない時$f(x)$ の符号は一定となり\cref{eq:2}はみたされず不適．

\paragraph{$f(x)=0$がただ一つの実数解を持つ時}

$f(x)=0$がただ一つの実数解$t$を持つ時，$f(x) = (x-t)^2 \ge 0$と書けるが，これは\cref{eq:2} を満たさず不適．

\paragraph{$f(x)=0$が二つの異なる実数解を持つ時}

$f(x)=0$が二つの異なる実数解を持つ時，$f(x) = (x-t)(x-s) \quad (t, s \in \mathbb{R}, t\neq s)$ と書ける．
「$t, s$ の少なくとも一方が $x \le 0, a \le x$ にある」と仮定する．
$t,s$の対称性からこれが $t$ だとして良い．この時 $g(x) = x - s$ とすると
\begin{align}
f(x)g(x) = (x-t)(x-s)^2 
\end{align}
の符号は $[0, a]$ で一定であり題意の条件式をみたさず矛盾．
よって背理法より $0 < t, s < a$ となる．

以上3つの場合分けから，$f(x)=0$が二つの異なる実数解を$0<x<a$に持つことが示された．

\end{document}
